\chapter{Hierarchical Graph Models: HPol}
\label{ch:hpol}
\draftnote{Target length: about 6 pages. Sources: \cite{alkhorem2026graph} (\PA) and
\cite{steiner2018hardening} (\PAA).}

\section{The Formalism}
\label{sec:hpol-formalism}
\draftnote{Policies as (type, subject, action, object) tuples over subject, action and
object hierarchies; the policy set as a hierarchical graph in which policies are links
\cite{alkhorem2026graph}.}

\section{Hierarchy as Abstraction}
\label{sec:hpol-hierarchy}
\draftnote{What a policy asserted at a parent node means for its children, and why this is
the approach's notion of abstraction.}

\section{Conflict and Redundancy Detection}
\label{sec:hpol-detection}
\draftnote{The positive--positive, positive--negative, negative--positive and
negative--negative cases, and the parent/child reconciliation that makes the partial
cases hard. Note that only the Subject hierarchy is walked.}

\section{How the Model Is Validated}
\label{sec:hpol-validation}
\draftnote{Exhaustive enumeration of model--request--response scenarios against a Python
implementation. What that form of evidence does and does not establish.}

\section{HPol as a Design-Transformation Tool}
\label{sec:hpol-transformation}
\draftnote{The Mutillidae case study \cite{steiner2018hardening}: a highest-privilege model,
a transformation to a least-privilege model, and the corresponding source-code change. A
refinement step performed informally on real code. Tracing one formalism from 2018 to 2026
is itself a finding.}

\section{Assessment}
\label{sec:hpol-assessment}
\draftnote{Position on the ladder: decidable and fast, but context-blind, solver-free, and
silent on whether the model-to-implementation step preserves anything. Fill the HPol rows of
\cref{tab:models,tab:evidence}.}
